% Source printed pp. 414-416 / PDF pp. 422-424.
\section{良商}\label{src:8.1}
\begin{epigraph}
我想讨论的基本问题是: 如果簇或概形 $V$ 参数化射影 $n$ 维空间 $\mathbb P^n$ 中某一类型的全部结构所成的集合或函子, 那么 $\mathbb P^n$ 的自同构群 $\mathrm{PGL}(n)$ 就在 $V$ 上作用. 在什么条件下存在商或轨道空间 $V/\mathrm{PGL}(n)$? 换言之, 何时能为这些结构构造足够多的``射影不变量''?
\epigraphauthor{David Mumford \cite[第 526 页]{Mum63}}
\end{epigraph}
设 $G$ 是代数闭域 $\Bbbk$ 上的仿射代数群, 作用于 $\Bbbk$ 上有限型代数空间 $U$. 对于视为 $\Bbbk$ 上群概形的有限抽象群 $G$, 我们已在下列情形中证明几何商 $U/G$ (定义~\bookref{5.2.1}) 的存在性:
\begin{itemize}
\item 若 $G$ 在 $U$ 上自由作用 (即作用映射 $G\times U\to U\times U$ 是单态射), 则 $U/G:=[U/G]$ 是 $\Bbbk$ 上有限型代数空间 (推论~\bookref{4.1.14}). 非有限情形也成立: 若 $G$ 是代数群且作用自由, 则 $[U/G]$ 是代数叠 (命题~\bookref{7.3.10}), 且 $[U/G]\to[U/G]\times[U/G]$ 是单态射, 因而 $U/G:=[U/G]$ 是代数空间 (定理~\bookref{5.5.10}).
\item 若 $U=\Spec A$ 仿射, 则 $U/G:=\Spec A^G$ 是几何商 (定理~\bookref{5.2.3}). 此外, 若 $U$ 射影 (分别地, 拟射影, 拟仿射), 则商 $U/G$ 是射影 (分别地, 拟射影, 拟仿射) 的 $\Bbbk$-概形 (习题~\bookref{5.2.14}).
\item 若 $U$ 分离, 则 Keel--Mori 定理 (\bookref{5.4.6}) 保证 $U/G$ 是分离代数空间. 非有限情形也成立: 若 $G$ 是仿射代数群, 作用的稳定子有限且约化, 并且作用映射 $G\times U\to U\times U$ 固有, 则 $[U/G]$ 是分离 Deligne--Mumford 叠 (定理~\bookref{4.6.4}), 几何商的存在性遂由 Keel--Mori 定理得到.
\end{itemize}
GIT 研究 $G$ 线性约化\footnote{GIT 可在约化代数群作用这一更一般的情形下发展, 但为简单起见, 我们仅考虑线性约化群; 见注记~\bookref{7.5.4}.}但未必有限的情形. 它允许点 $u\in U$ 的稳定子 $G_u$ 不有限, 轨道 $Gu$ 不闭, 例如 $\mathbb G_m$ 在 $\mathbb A^1$ 上的作用. 我们已讨论仿射 GIT (\bookref{7.5.13}): 若 $G$ 作用于仿射 $\Bbbk$-概形 $\Spec A$, 则有交换图
\[
\begin{tikzcd}[ampersand replacement=\&]
\Spec A\arrow[d]\arrow[dr,"\widetilde\pi"]\&{}\\
{[\Spec A/G]}\arrow[r,"\pi"']\&{(\Spec A)\sslash G:=\Spec A^G}.
\end{tikzcd}
\]
其中 $\pi:[\Spec A/G]\to\Spec A^G$ 是良模空间. 映射 $\widetilde\pi:\Spec A\to\Spec A^G$ 是范畴商, 且两个 $\Bbbk$-点的像相同当且仅当它们的轨道闭包相交; 我们称此映射为良商.
\subsection{定义与初步性质}\label{src:8.1.1-subsection}
\begin{definition}[良商]\label{src:8.1.1}
设代数闭域 $\Bbbk$ 上的线性约化群 $G$ 作用于 $\Bbbk$ 上代数空间 $U$. 若 $G$-不变映射 $\widetilde\pi:U\to X$ 满足下列条件, 则称之为\newterm{良商}{good quotient}:\footnote{\cite[定义 1.5]{Ses72} 对任意仿射代数群 $G$ 将良商定义为仿射 $G$-不变态射 $\widetilde\pi:U\to X$, 满足 $\mathcal O_X\xrightarrow{\sim}(\pi_*\mathcal O_U)^G$, 以及命题~\bookref{8.1.3}(1)--(2) 的性质. 当 $G$ 线性约化时, 两种定义相同.}
\begin{enumerate}[(1)]
\item $\mathcal O_X\to(\pi_*\mathcal O_U)^G$ 是同构 (对平展 $X$-概形 $V$, 有 $(\pi_*\mathcal O_U)^G(V)=\Gamma(U_V,\mathcal O_{U_V})^G$);
\item $\widetilde\pi$ 仿射.
\end{enumerate}
$U$ 对 $G$ 的良商常记为 $U\sslash G=X$.
\end{definition}
\begin{remark}\label{src:8.1.2}
映射 $\widetilde\pi:U\to X$ 是良商, 当且仅当 $\pi:[U/G]\to X$ 是良模空间. 因为这两种性质均为平展局部性质 (引理~\bookref{7.5.24}(1)), 验证等价性时可设 $X=\Spec B$ 仿射. 对 $(\Rightarrow)$, $U=\Spec A$ 也仿射且 $B=A^G$, 因而 $[\Spec A/G]\to\Spec A^G$ 是良模空间. 对 $(\Leftarrow)$, 注意 $U\to[U/G]$ 仿射, 且 $\pi_*$ 在拟凝聚层上正合, 所以推前 $\widetilde\pi_*$ 在拟凝聚层上正合; 由 Serre 仿射性判据 (\bookref{5.5.16}) 得 $\widetilde\pi$ 仿射.
\end{remark}
\begin{proposition}\label{src:8.1.3}
设代数闭域 $\Bbbk$ 上的线性约化群 $G$ 作用于 $\Bbbk$ 上代数空间 $U$. 若 $\widetilde\pi:U\to X$ 是良商, 则:
\begin{enumerate}[(1)]
\item $\widetilde\pi$ 满射, 且闭 $G$-不变子概形的像闭. 对来自概形的态射 $T\to X$ 作基变换后, 结论仍成立;
\item 对闭 $G$-不变子概形 $Z_1,Z_2\subseteq U$, 有 $\im(Z_1\cap Z_2)=\im(Z_1)\cap\im(Z_2)$. 特别地, 对 $u_1,u_2\in U(\Bbbk)$, $\widetilde\pi(u_1)=\widetilde\pi(u_2)$ 当且仅当 $\overline{Gu_1}\cap\overline{Gu_2}\ne\emptyset$; $\widetilde\pi$ 诱导 $U$ 中闭 $G$-轨道与 $X$ 的 $\Bbbk$-点之间的双射;
\item 若 $U$ 是 Noether 的, 则 $X$ 也是. 若 $U$ 在 $\Bbbk$ 上有限型, 则 $X$ 也是; 对每个带 $G$-作用的凝聚 $\mathcal O_U$-模 $F$, $(\pi_*F)^G$ 凝聚;
\item $\widetilde\pi$ 对到代数空间的 $G$-不变映射是万有的.
\end{enumerate}
\end{proposition}
\begin{proof}
因为 $[U/G]\to X$ 是良模空间, 这由定理~\bookref{7.5.11} 得到.
\end{proof}
\begin{remark}[GIT 中的半稳定约化]\label{src:8.1.4}
因为 $[U/G]\to X$ 普遍闭 (定理~\bookref{7.5.11}(1)), 它满足普遍闭性的赋值判别法 (定理~\bookref{4.8.7}). 这意味着: 对每个 $\Bbbk$ 上分式域为 $K$ 的 DVR $R$, 以及有提升 $\eta:\Spec K\to U$ 的映射 $\Spec R\to X$, 存在 DVR 扩张 $R\to R'$, $R'$ 的分式域 $K'$ 上的元素 $g'\in G(K')$, 以及下列交换图中的提升:
\[
\begin{tikzcd}[ampersand replacement=\&,column sep=large]
\Spec K'\arrow[r]\arrow[d]\arrow[rr,bend left=28,"{g'\cdot\eta|_{K'}}"]\&\Spec K\arrow[r,"\eta"]\arrow[d]\&U\arrow[d,"\widetilde\pi"]\\
\Spec R'\arrow[r]\arrow[urr,dashed]\&\Spec R\arrow[r]\&X.
\end{tikzcd}
\]
事实上, 若 $R=\Bbbk[\![x]\!]$, 可令 $R\to R'$ 有限; 见 \cite[引理 5.3]{Mum77} 与 \cite[定理 A.8]{AHLH23}.
\end{remark}
